\section{动作表示、API 边界与架构可塑性}

复杂环境只有在动作可表达、可组合时才可学习。本节从 AlphaStar 的结构化 action head 过渡到 LLM tool call，关注模型到底能选择哪些动作、参数如何分解，以及执行副作用怎样被系统约束；这些接口决定后续 RL 能探索到什么策略。

在 17--18 分钟段，讲者把 AlphaStar 的 action head 与今天 LLM 的 Tool Use 做了直接映射。这是全讲最有工程价值的部分之一：它揭示了``动作接口设计''本身就是能力上限。

\begin{figure}[H]
\centering
\includegraphics[width=0.92\textwidth]{frame-02.jpg}
\caption{AlphaStar 的 autoregressive action head}
\label{fig:action-head}
\end{figure}

\begin{knowledgebox}{读图：Autoregressive head 如何缩小动作搜索}
模型不是一次从 $10^{26}$ 个组合动作中硬选一个，而是依次预测动作类型、延迟、队列、单位集合和目标位置，每个子头都条件于前面的选择。LLM function calling 同样应把工具、实体、参数和权限拆开；若所有内容塞进自由文本，系统就失去类型检查与局部纠错。
\end{knowledgebox}
图 \ref{fig:action-head} 展示了动作拆解的结构\footnote{来源：\href{https://www.youtube.com/watch?v=ntjOxjZMaac}{Lecture 03 视频帧}，时间戳 00:17:40。}：What（动作类型）、When（时机/延迟）、Who（施加对象）、Where（目标位置）。这与 LLM function calling 中 ``tool name + arguments + execution context'' 基本同构。

\begin{knowledgebox}{从 StarCraft Action 到 LLM Tool Call 的映射}
\begin{itemize}
    \item What $\leftrightarrow$ 选择工具（search / code / browser / retrieval）
    \item When $\leftrightarrow$ 何时调用（先规划后执行，还是边思考边调用）
    \item Who $\leftrightarrow$ 对哪个实体施加动作（文件、进程、用户会话、数据库记录）
    \item Where $\leftrightarrow$ 操作落点（参数空间、页面节点、API endpoint）
\end{itemize}
\end{knowledgebox}

\begin{importantbox}{接口比模型更先成为瓶颈}
如果 API 的动作语义过粗，模型无法表达精细计划；如果 API 过细，探索空间会指数增长。课程给出的工程建议是：先建立``可组合原语''，再让策略在原语之上学习分层调用。
\end{importantbox}

\begin{warningbox}{常见误区：把 Tool Use 仅当作检索插件}
把 tool 视为``补充知识''会忽视其控制能力（actuation）。在 Agent 任务里，tool 既是信息通道，也是行为通道。行为通道失控将直接导致越权、幻觉执行、不可逆副作用。
\end{warningbox}

为了使该思想可落地，可以把动作接口写成显式的约束语法：
\begin{lstlisting}[caption={一个最小化的 Tool API Schema（示意）}]
TOOL_CALL := {
  "tool": one_of["search", "terminal", "editor", "web_click"],
  "goal": natural_language_intent,
  "args": structured_arguments,
  "safety_level": one_of["read_only", "bounded_write", "privileged"],
  "rollback_plan": optional_steps
}
\end{lstlisting}
这类 schema 的价值在于：它让训练目标从``生成像人类的文本''转成``生成可验证的动作计划''。

\subsection{本章小结}
动作空间定义了可学能力边界。AlphaStar 的经验告诉我们，LLM Agent 的架构设计应优先解决动作可分解、可约束、可验证三个问题，否则后续 RL 和多 Agent 训练只是在不稳定接口上叠加复杂度。

